¶ 11) 对任意整数 \(n \geq 1\) ，定义映射 \(\varphi_n\)，其定义由 \((0, \frac{1}{2} (\times \mathbf{T}^2 \text{ 到 } \mathbf{R}^3 \text{ 令 } \varphi_n(r; \alpha, \beta) = ((n + r \cos 2\pi\beta) \cos 2\pi\alpha, (n + r \cos 2\pi\beta) \sin 2\pi\alpha, r \sin 2\pi\beta)\) 给出。令

\[
\mathrm{T} _ {n} = \varphi_ {n} \left(\left[ 0, \frac {1}{2} \right[ \times \mathbf {T} ^ {2}\right), \quad \mathrm{S} _ {n} = \varphi_ {n} (\{0 \} \times \mathbf {T} ^ {2}).
\]

\begin{enumerate}
\def\labelenumi{\alph{enumi})}
\item
  证明 \(\varphi_{n}\) 在 \(\left[0,\frac{1}{2}\right[\times T^{2}\) 上的限制是一个（类为 \(C^{\omega}\) 的）从 \(\left]0,\frac{1}{2}\right[\times T^{2}\) 到 \(T_{n}-S_{n}\) 的同构。
\item
  设 \(f_{n}:\mathrm{T}_{n}\to \mathrm{T}_{n}\) 为在 \(\mathrm{S}_n\) 上与恒等映射重合并满足下式的映射：
\end{enumerate}

\[
f _ {n} (\varphi_ {n} (r; \alpha , \beta)) = \varphi_ {n} (r; \alpha - (n - 1) \beta , - \alpha + n \beta)
\]

当 r \textgreater{} 0 时，证明 \(f_{n}\) 诱导出 \(T_{n} - S_{n}\) 的一个自同构。

\begin{enumerate}
\def\labelenumi{\alph{enumi})}
\setcounter{enumi}{2}
\tightlist
\item
  证明在 \(R^{3}\) 上存在一个实解析流形结构，使得映射 \(f_{n}: T_{n} \to R^{3}\) 和典则注入 \(R^{3}-\bigcup_{n} S_{n} \to R^{3}\) 都是解析的。将由此定义的实解析流形记为 X。
\end{enumerate}

\(d)\) 对于 \(\vartheta \in \mathbf{T}\)，将旋转记为 \(\mathrm{R}_{\vartheta}\)

\[
(x, y, z) \mapsto (x \cos 2 \pi \vartheta - y \sin 2 \pi \vartheta , x \sin 2 \pi \vartheta + y \cos 2 \pi \vartheta , z)
\]

在 \(\mathbf{R}^3\) 上。证明，对于 \(\theta \in \mathbf{T}\) 和 \(u \in \mathrm{X}\)，规定

\[
\vartheta . u = \mathrm{R} _ {\vartheta} (u) \quad \text { if } \quad u \in \mathrm{X} - \bigcup_ {n} \mathrm{S} _ {n};
\]

\(\vartheta .u = \mathrm{R}_{n\vartheta}(u)\) 若 \(u\in \mathrm{S}_n\) 对某个整数 \(n\geq 1\)

定义了 \(\mathbf{T}\) 在 \(\mathrm{X}\) 上的解析作用律。

\begin{enumerate}
\def\labelenumi{\alph{enumi})}
\setcounter{enumi}{4}
\item
  证明 X 的轨道类型集合（对于 d) 中定义的 T 作用）是无限的。
\item
  证明不存在一个与 T 的作用相容的嵌入，将 X 嵌入某个 T 在线性作用的有限维向量空间中（使用习题 9）。
\end{enumerate}

\begin{enumerate}
\def\labelenumi{\arabic{enumi})}
\setcounter{enumi}{11}
\tightlist
\item
  设 \(\mathrm{G}\) 为紧李群。
\end{enumerate}

\begin{enumerate}
\def\labelenumi{\alph{enumi})}
\item
  证明整子群的正规化子的共轭类集合是有限的（考察 G 在 L(G) 的子空间 Grassmann 流形上的作用，并应用习题 9）。
\item
  （紧）半单子群的共轭类集合是有限的（注意到一个李代数至多包含有限多个半单理想（第一章，§6，习题 7），并使用 a)）。
\end{enumerate}

\begin{enumerate}
\def\labelenumi{\arabic{enumi})}
\setcounter{enumi}{12}
\item
  设 G 为李群，H、K 为 G 的两个紧子群。假设 G 具有忠实的有限维线性表示。证明，存在 H 的子群的有限集合 F，使得对于所有 \(g \in G\)，子群 \(H \cap gKg^{-1}\) 在 H 中共轭于 F 中的某个子群（使用习题 5 和 9）。
\item
  假设流形 X 是仿紧的且局部有限维。设 G 为在 X 上真作用的李群，H 为 G 的紧子群，t 为 H 的共轭类。
\end{enumerate}

\begin{enumerate}
\def\labelenumi{\alph{enumi})}
\item
  证明 X 中固定子群等于 H 的点所成的集合 \(X_{H}\) 是 X 的局部闭子流形。
\item
  证明映射 \((g,x)\mapsto gx\) 从 \(G\times X_{H}\) 到 X 诱导出一个（类为 \(C^{r}\) 的）从 \(G\times^{\mathrm{N(H)}}X_{\mathrm{H}}\) 到 \(\mathrm{X}_{(t)}\) 的同构。
\end{enumerate}

¶ 15) 设 G 为在拓扑空间 E 上真作用的局部紧拓扑群；设 \(\rho\) 是 G 在有限维实向量空间 V 上的一个表示。记 G 上的右 Haar 测度为 dg。

\begin{enumerate}
\def\labelenumi{\alph{enumi})}
\tightlist
\item
  设 \(\mathcal{P}\) 为 E 的具有如下性质的子集 A 所成的集合：存在 E 的一个开覆盖 \((\mathrm{U}_{\alpha})_{\alpha \in \mathrm{I}}\)，使得对于所有 \(\alpha \in \mathrm{I}\)，\(g \in \mathrm{G}\) 中满足 \(g\mathrm{A} \cap \mathrm{U}_{\alpha} \neq \emptyset\) 的元素所成的集合在 G 中相对紧。
\end{enumerate}

证明，对于每个支集属于 P 的连续函数 \(f: E \to V\)，映射 \(x \mapsto \int_{G} \rho(g)^{-1} f(gx) dg\) 是从 E 到 V 的连续映射，并且与 G 的作用相容。

\begin{enumerate}
\def\labelenumi{\alph{enumi})}
\setcounter{enumi}{1}
\tightlist
\item
  作如下两种假设之一：
\end{enumerate}

\begin{enumerate}
\def\labelenumi{(\roman{enumi})}
\item
  空间 E/G 是正则的；
\item
  存在 E 的一个开子集 U 和 G 的一个紧子集 K，使得 \(\mathrm{E} = \mathrm{GU}\)，并且 \(g\mathrm{U}\cap \mathrm{U} = \varnothing\) 对 \(g\notin \mathrm{K}\) 成立。
\end{enumerate}

证明 E 中每一点 x 都有一个属于 P 的邻域（在情形 (i) 中，取 \(A = V \cap W\)，其中 V 是 x 的一个邻域，当 g 在 G 的某个紧子集之外时满足 \(gV \cap V = \varnothing\)；W 取为 Gx 的一个在 G 下稳定且包含于 GV 中的闭邻域）。

E 中每一点都有一个在 G 下稳定且满足 (ii) 的开邻域。

\begin{enumerate}
\def\labelenumi{\alph{enumi})}
\setcounter{enumi}{2}
\tightlist
\item
  另外假设 E 是完全正则的。设 \(x \in E, v \in V\)，使得 x 的固定子群包含于 v 的固定子群。证明存在连续映射 \(F : E \to V\)，它与 G 的作用相容，且 \(\mathrm{F}(x) = v\)。（令 F 为 E 上支集属于 P 的连续数值函数空间，并令 \(u : F \to V\) 为映射 \(\alpha \mapsto \int_{G} \alpha(gx) \rho(g^{-1}). v \, dg\)。设 C 是 V 中 v 的一个凸邻域；构造一个属于 P 的 x 的邻域 A，使得 \(gx \in A\) 蕴含 \(\rho(g^{-1}). v \in C\)，并构造 E 上的一个函数 \(\alpha\)，其支集包含于 A，满足 \(\alpha(x) \neq 0\) 且 \(\int_{G} \alpha(gx) \, dg = 1\)。证明 \(u(\alpha)\) 属于 C，并推出 \(v \in Im u\)。）
\end{enumerate}

\begin{enumerate}
\def\labelenumi{\arabic{enumi})}
\setcounter{enumi}{15}
\tightlist
\item
  设 G 为在分离拓扑空间 E 上真作用的拓扑群。设 x 是 E 中一点，H 是它在 G 中的固定子群，S 是包含 x 且在 H 下稳定的 E 的子集。群 H 右作用于 \(G \times S\)，作用公式为 \((g, s).h = (gh, h^{-1}s)\)，其中 \(g \in G, h \in H, s \in S\)；映射 \((g, s) \mapsto gs\) 通过取商诱导出从 \((\mathrm{G} \times \mathrm{S})/\mathrm{H}\) 到 X 的映射。若此映射是到 X 的某个开子集上的同胚，则称 S 是 x 处的横截集。
\end{enumerate}

\begin{enumerate}
\def\labelenumi{\alph{enumi})}
\item
  证明，若 \(\mathbf{G}\) 是离散的，则 \(x\) 处存在一个横截集。
\item
  设 F 是一个分离拓扑空间，G 在其上真作用；设 \(f : E \to F\) 是与 G 的作用相容的连续映射，且 \(f(x)\) 在 G 中的固定子群等于 H。证明，若 S 是 \(f(x)\) 处的横截集，则 \(f^{-1}(S)\) 是 x 处的横截集。
\item
  设 N 是 G 的闭正规子群，\(\pi: E \to E/N\) 是典则投影，T 是 E/N 上 G/N 作用下 \(\pi(x)\) 处的横截集，且 \(S \subset \pi^{-1}(T)\) 是 \(\pi^{-1}(T)\) 上 HN 作用下 x 处的横截集。证明 S 是 E 中 x 处的横截集（对于 G 的作用）。
\end{enumerate}

\begin{enumerate}
\def\labelenumi{\arabic{enumi})}
\setcounter{enumi}{16}
\tightlist
\item
  设 \(G\) 为李群。我们将证明 \(G\) 具有如下性质：
\end{enumerate}

\begin{enumerate}
\def\labelenumi{(\Alph{enumi})}
\setcounter{enumi}{19}
\tightlist
\item
  对于每个完全正则拓扑空间 \(\mathrm{E}\)，在其上 \(\mathrm{G}\) 真作用，并且对于每一点 \(x\)，在 \(\mathrm{E}\) 中的 \(x\) 处都存在一个横截集（习题 16）。
\end{enumerate}

\begin{enumerate}
\def\labelenumi{\alph{enumi})}
\item
  证明每个具有忠实有限维线性表示的李群都满足 (T)（使用习题 15 和 16 b)，以及命题 6）。
\item
  证明，若 G 有一个闭正规子群 N，使得 G/N 满足 (T)，并且对 G 的每个紧子群 K，KN 满足 (T)，则 G 满足 (T)（应用习题 16 c)）。
\item
  若 \(\mathrm{G}_0\) 是紧的，则 G 满足 (T)。
\item
  证明，若 G 有一个离散正规子群 N，使得 G/N 满足 (T)，则 G 满足 (T)。
\item
  证明 G 满足 (T)（令 N 为伴随表示的核；证明 \(G/N_{0}\) 满足 (T)，然后应用 a)、b) 以及 §1 的习题 9）。
\end{enumerate}

\begin{enumerate}
\def\labelenumi{\arabic{enumi})}
\setcounter{enumi}{17}
\tightlist
\item
  设 \(G\) 为在完全正则拓扑空间 \(E\) 上真作用的李群。
\end{enumerate}

\begin{enumerate}
\def\labelenumi{\alph{enumi})}
\item
  设 \(x \in E\)，且 t 是它的轨道类型；证明存在 x 的一个在 G 下稳定的开邻域 U，使得对于所有 \(u \in U\)，u 的类型都满足 \(\geq t\)。
\item
  假设 G 在 E 上自由作用；令 \(\pi: E \to E/G\) 为典则投影。证明，对于 E/G 中的每一点 z，都存在 z 的一个开邻域 U 和连续映射 \(s: U \to E\)，使得 \(\pi \circ s(u) = u\) 对所有 \(u \in U\) 都成立。
\end{enumerate}

\begin{enumerate}
\def\labelenumi{\arabic{enumi})}
\setcounter{enumi}{18}
\tightlist
\item
  设 G 为李群，H 为 G 的紧子群。证明存在 H 的一个邻域 V，使得 G 中包含于 V 的每个子群都共轭于 H 的某个子群（将习题 18 a) 应用于 G 的紧子集空间，参见~《积分》，第八章，§5，第 6 节）。
\end{enumerate}

¶ 20) 设 G 为紧李群，m 为正整数。

\begin{enumerate}
\def\labelenumi{\alph{enumi})}
\item
  证明 G 中满足阶 \(\leq m\) 的子群的共轭类集合是有限的（假设结论不成立，构造一个有限群 F 以及一列同态 \(\varphi_{n}: F \to G\)，使得 \(\varphi_{i}(F)\) 不共轭于 \(\varphi_{j}(F)\)（当 \(i \neq j\) 时），并且 \(\varphi_{n}(f)\) 都趋于极限 \(\varphi(f)\)，对所有 \(f \in F\) 均如此；证明 \(\varphi\) 是同态，并证明这将导致与习题 19 的矛盾）。
\item
  证明 G 中所有元素的阶均满足 \(\leq m\) 的子群 F 的共轭类集合是有限的（令 \(\mu\) 为 G 上的 Haar 测度，令 U 为单位元的一个对称邻域，使得 \(U^{2}\) 不包含阶满足 \(\leq m\) 的非平凡元素；证明 \(\operatorname{Card}(F) \leq \mu(G)/\mu(U)\)）。
\end{enumerate}

¶ 21) 设 G 为紧李群，T 为 G 的极大环面。

\begin{enumerate}
\def\labelenumi{\alph{enumi})}
\item
  设 S 是 G 的闭子群的一个在共轭下稳定的集合，并且子群族 \((\mathrm{S} \cap \mathrm{T})_{\mathrm{S} \in \mathcal{S}}\) 是有限的。证明子群 \(S_{0}\)（其中 \(S \in S\)）的共轭类集合是有限的（利用习题 12 a)），化归到子群 \(S_{0}\) 都是正规的情形；考察群 \(\mathrm{C}(S_{0})_{0}\) 和 \(\mathrm{D}(S_{0})\)，并应用习题 12 b)）。
\item
  证明 S 是 G 的子群的有限个共轭类的并（利用 a)，化归到子群 \(S \in S\) 都具有同一个单位元连通分支 \(\Sigma\) 的情形，其中该分支在 G 中正规；然后给出群 \(S/\Sigma\) 中元素的阶的上界，并应用习题 20 b)）。
\item
  设 E 是一个 G 在其上连续作用的分离拓扑空间。证明，若 E 的点关于 T 的作用只有有限种轨道类型，则关于 G 的作用也只有有限种轨道类型。
\end{enumerate}

\exercisesection{附录~I}

\begin{enumerate}
\def\labelenumi{\arabic{enumi})}
\tightlist
\item
  设 \(G\) 为连通紧群。以 \(d(G)\) 表示 \(G\) 的李群商群的维数上界。假设 \(d(G) < \infty\)。
\end{enumerate}

\begin{enumerate}
\def\labelenumi{\alph{enumi})}
\item
  设 K 是 G 的闭正规子群；证明 \(d(\mathrm{G}/\mathrm{K}) \leq d(\mathrm{G})\)，并且当 K 全不连通时 \(d(\mathrm{G}/\mathrm{K}) = d(\mathrm{G})\)。
\item
  证明 \(\mathrm{D}(\mathbf{G})\) 是李群，并且同态 \((x,y)\mapsto xy\) 从 \(\mathrm{C}(\mathrm{G})_0\times \mathrm{D}(\mathrm{G})\) 到 \(\mathbf{G}\) 的核是有限的。
\item
  设 \(p = d(\mathrm{C}(\mathrm{G})_{0})\)。则 \(p < \infty\)；证明存在一个全不连通的紧群 D 以及一个具有稠密像的同态 \(i : Z^{p} \to D\)，使得 \(\mathrm{C}(\mathrm{G})_{0}\) 同构于 \((\mathbf{R}^{p} \times \mathrm{D}) / \Gamma\)，其中 \(\Gamma\) 是 \(Z^{p}\) 在同态 \(x \mapsto (x, i(x))\) 下的像（将 \(\mathrm{C}(\mathrm{G})_{0}\) 写成维数为 p 的环面的射影极限）。
\item
  假设 G 是局部连通的；证明 G 从而是李群。
\end{enumerate}

\specialchapter{记号索引}

这些参考编号分别表示章、节和编号。

\[
k \colon \mathrm{VII.Conventions}.
\]

\[
\mathrm{V} _ {\lambda} (\mathrm{S}), \mathrm{V} _ {\lambda} (s), \mathrm{V} ^ {\lambda} (\mathrm{S}), \mathrm{V} ^ {\lambda} (s) \colon \text { VII.1.1 }.
\]

\[
\mathfrak {g} _ {\lambda} (\mathfrak {h}), \mathfrak {g} ^ {\lambda} (\mathfrak {h}) \colon \text { VII.1.3 }.
\]

\[
a _ {i} (x), \operatorname{rk} (\mathfrak {g}): \text { VII.2.2 }.
\]

\[
\operatorname{Aut} (\mathfrak {g}), \operatorname{Aut} _ {e} (\mathfrak {g}): \text { VII.3.1 }.
\]

\[
\mathscr {C} ^ {\infty} \mathfrak {g} = \cap \mathscr {C} ^ {n} \mathfrak {g}: \text { VII.3.4. }
\]

\[
r _ {\rho}, r _ {\rho} ^ {0}: \text { VII.4.1 }.
\]

\[
e (\mathfrak {g}): \text {   VII.5.2.   }
\]

\[
\mathrm{A} _ {\mathrm{V}} \text {: VII.App.I.1. }
\]

\[
\mathrm{D} f \text {: VII.App.I.2.}
\]

\[
H, X _ {+}, X _ {-}: \text {   VIII.1.1.   }
\]

\[
\mathrm{V} (m): \text { VIII.1.3 }.
\]

\[
h (t), \theta (t): \text {   VIII.1.5.   }
\]

\[
\mathfrak {g}, \mathfrak {h}, \mathrm{R} (\mathfrak {g}, \mathfrak {h}) = \mathrm{R}: \text {   VIII.2.2.   }
\]

\[
\mathfrak {g} ^ {\alpha}, \mathfrak {h} _ {\alpha}, \mathfrak {s} _ {\alpha}, H _ {\alpha}, h _ {\alpha}, \theta_ {\alpha} (t), s _ {\alpha}: \text {   VIII.2.2.   }
\]

\[
\mathfrak {h} _ {\mathbf {Q}}, \mathfrak {h} _ {\mathbf {Q}} ^ {*}, \mathfrak {h} _ {\mathbf {R}}, \mathfrak {h} _ {\mathbf {R}} ^ {*}: \text {   VIII.2.2.   }
\]

\[
\mathrm{N} _ {\alpha \beta}: \mathrm{VIII.2.4}.
\]

\[
\mathfrak {g} ^ {\mathrm{P}}, \mathfrak {h} _ {\mathrm{P}}: \text {   VIII.3.1.   }
\]

\[
n (\alpha , \beta), X _ {\alpha} ^ {0}, H _ {\alpha} ^ {0}, X _ {- \alpha} ^ {0}: \text { VIII.4.2 }.
\]

\[
\theta , \mathfrak {a}, \mathfrak {a} _ {+}, \mathfrak {a} _ {-}: \text {   VIII.4.2.   }
\]

\[
x _ {\alpha \beta}, y _ {\alpha \beta}, \mathfrak {n}: \text {   VIII.4.3.   }
\]

\[
\operatorname{Aut} (\mathfrak {g}, \mathfrak {h}), \mathrm{A} (\mathrm{R}): \text {   VIII.5.1.   }
\]

\[
\mathrm{T} _ {\mathrm{P}}, \mathrm{T} _ {\mathrm{Q}}, \operatorname{Aut} _ {0} (\mathfrak {g}), \operatorname{Aut} _ {0} (\mathfrak {g}, \mathfrak {h}), \operatorname{Aut} _ {e} (\mathfrak {g}, \mathfrak {h}): \text {VIII.5.2}.
\]

\[
(\mathfrak {g}, \mathfrak {h}), \check {\mathrm{R}}, \mathrm{W}, \mathrm{B}, \mathrm{R} _ {+}, \mathrm{R} _ {-}, \mathfrak {n} _ {+}, \mathfrak {n} _ {-}, \mathfrak {b} _ {+}, \mathfrak {b} _ {-}, X _ {\alpha}, Y _ {\alpha}, H _ {\alpha}: \text {VIII.6.Conventions.}
\]

\[
w _ {0} \colon \text {   VIII.7.2.   }
\]

\[
[ \mathrm{E} ], [ \mathrm{E} ] ^ {*}, \mathscr {R} (\mathfrak {a}): \text {   VIII.7.6.   }
\]

\[
\mathbf {Z} [ \Delta ], e ^ {\lambda}, \operatorname{ch} (\mathrm{E}) \colon \text { VIII.7.7 }.
\]

\[
\exp (\lambda): \text {   VIII.8.1.   }
\]

\[
\theta (a), \theta^ {*} (a), \mathrm{I} (\mathfrak {g}), \mathrm{I} (\mathfrak {g} ^ {*}), \mathrm{S} (\mathfrak {g}), \mathrm{S} (\mathfrak {g} ^ {*}) \colon \text { VIII.8.3 }.
\]

\[
u ^ {\sharp}, \chi_ {\lambda}, \eta , \delta , \theta , \varphi : \mathrm{VIII.8.5}.
\]

\[
\mathfrak {B} (\nu), \mathrm{K}, d: \text {   VIII.9.1.   }
\]

\[
\mathrm{J} (e ^ {\mu}) \colon \text { VIII.9.2. }
\]

\((x,h,y)\) : VIII.11.1. \(x^{(n,d)},x^{(n)},\binom{x}{n}\) : VIII.12.2. \(\mathbf{x}^{(n)}\) , \([n],c_r,\mathcal{G}\) : VIII.12.6. \(\mathcal{H},\mathcal{U}_{+},\mathcal{U}_{-},\mathcal{U}_{0},\mathcal{L},\mathcal{G}\) : VIII.12.6. \(\mathfrak{o}(\Psi)\) : VIII.13.2. \(\mathfrak{sp}(\Psi)\) : VIII.13.3. \(\mathrm{G}_0,\mathrm{C}(\mathrm{G}),\mathrm{D}(\mathrm{G})\) : IX.1.Conventions. \(\mathrm{N_G(H),Z_G(H),N(H),Z(H)}\) : IX.1.Conventions. \(\Gamma (\mathrm{G})\) : IX.1.2. \(\tilde{\mathrm{D}} (\mathrm{G})\) : IX.2.3.\\
rk G: IX.2.4. \(\mathrm{W_G(T),W(T)}\) : IX.2.5. \(\mathfrak{a}_{[\mathbb{R}]},\mathfrak{g}_{(\mathbb{C})},\mathfrak{g}_{\mathbb{C}}\) : IX.3.1. \(\mathfrak{a}_u,u_\alpha ,v_\alpha\) : IX.3.2. \(\mathbf{U}(\varPhi),\mathbf{SU}(\varPhi),\mathfrak{u}(\varPhi),\mathfrak{s}\mathfrak{u}(\varPhi)\) : IX.3.4. \(\mathbf{U}(n,\mathbf{C}),\mathbf{SU}(n,\mathbf{C}),\mathfrak{u}(n,\mathbf{C}),\mathfrak{s}\mathfrak{u}(n,\mathbf{C})\) : IX.3.4. \(\mathbf{O}(\mathrm{Q}),\mathbf{SO}(\mathrm{Q}),\mathfrak{o}(\mathrm{Q})\) : IX.3.5. \(\mathbf{O}(n,\mathbf{R}),\mathbf{SO}(n,\mathbf{R}),\mathfrak{o}(n,\mathbf{R})\) : IX.3.5. \(U,V,K,\theta\) : IX.3.6.\\
X(H),X(f): IX.4.1. \(\delta (a)\) : IX.4.1. \(\Gamma (f)\) : IX.4.2. \(\langle a,X\rangle\) : IX.4.2. \(\tilde{\rho},\tilde{\mathrm{V}},\tilde{\mathrm{V}}_{\lambda}(\mathrm{G})\) : IX.4.3.\\
P(ρ,T): IX.4.3. \(\tilde{\mathrm{L}} (\rho)\) : IX.4.3.\\
R(G,T): IX.4.4.\\
\(\mathrm{V}(\alpha),\mathrm{K}(\alpha),\iota_{\alpha},\rho_{\alpha}\): IX.4.4.\\
\(\mathrm{Z}_{\alpha},\mathrm{S}_{\alpha}\): IX.4.5.\\
\(\mathrm{K}_{\alpha}\): IX.4.5.\\
\(\mathrm{R}^{\vee}(\mathrm{G},\mathrm{T})\): IX.4.5.\\
N(G,T): IX.4.6.\\
D*(G,T),D*(G,T): IX.4.9.\\
D*(f),D*(f): IX.4.9.\\
D*(G),D*(G): IX.4.9.\\
Aut(G),Aut(G,T): IX.4.10.\\
\(\mathrm{G}_{r},\mathrm{T}_{r}\): IX.5.1.\\
\(\mathrm{W}_{\alpha},\mathrm{W}'_{\alpha}\): IX.5.2.\\
\(\mathrm{H}_{\alpha,n},t_{r}\): IX.5.2.\\
\(\mathrm{H}_{A}\): IX.5.2.\\
\(\mathfrak{g}_{\mathrm{reg}}\): IX.5.2.\\
\(f,f_{r}\): IX.5.4.\\
\(\varphi,\varphi_{r},\varphi_{A}\): IX.5.4.\\
\(\mathfrak{g},r,\psi_{r}\): IX.5.4

\(w(\mathrm{G})\) : IX.6.Conventions. \(u\cap v\) : IX.6.1. \(\omega_{\mathrm{G / T}}\cap \omega_{\mathrm{T}}\) : IX.6.2. \(\mathrm{Ad}_{\mathfrak{g} / \mathfrak{t}}(t),\delta_{\mathrm{G}}(t)\) : IX.6.2. \(\mathrm{R}_+(\mathrm{G},\mathrm{T})\) : IX.6.2. \(\lambda_{\mathfrak{h}}\) : IX.6.3. \(\pi_{\mathfrak{g}}\) : IX.6.3. \(\mathcal{S}^r (\mathrm{X};\mathrm{E})\) : IX.6.4. \(s^\sharp\) : IX.6.4. \(\Omega (\mathrm{X}),\Omega (\mathrm{X})^{\mathrm{G}}\) : IX.6.5. \(\Gamma (\mathrm{T})_{++},\mathrm{X}_{+},\mathrm{R}_{+},\mathrm{R}_{-}\) : IX.7.Conventions. \(\mathrm{X}_{++}\) : IX.7.1. \(\varpi_{\alpha},\rho\) : IX.7.1. \(\Theta (\mathrm{G})\) : IX.7.2. \(\operatorname {R}(\operatorname {G})\) : IX.7.3. \([\tau ],\operatorname {Ch}(\tau)\) : IX.7.3. \(\mathrm{Z}\Theta (\mathrm{G})\) : IX.7.3. \(^w f,\mathrm{J}(f)\) : IX.7.4. \(\chi_{\lambda}\) : IX.7.4. \(\mathrm{ZL}^2 (\mathrm{G})\) : IX.7.4. \(\mathrm{Q}_{\Gamma},\tilde{\Gamma} (\tau)\) : IX.7.6. \(\hat{\mathbf{G}}\) : IX.8.1. \(\operatorname {E}_u,d(u),\| \mathrm{A}\|_\infty ,\| \mathrm{A}\| _2,\langle \mathrm{A}|\mathrm{B}\rangle\) : IX.8.1.\\
F( \(\hat{\mathrm{G}}\) ), \(\mathrm{L}^2 (\mathrm{G})\) : IX.8.1. \(u(f),\overline{\mathcal{F}} (f)\) : IX.8.1. \(\mathcal{F}_u\mathrm{A},\mathcal{F}(\mathrm{A})\) : IX.8.1. \(\gamma (s)f,\delta (s)f\) : IX.8.1. \(\mathrm{L}_t f,\mathrm{R}_t f\) : IX.8.2. \(\lambda (u),\tilde{\Gamma} (u)\) : IX.8.2. \(\varphi \preccurlyeq \psi\) : IX.8.2. \(\mathcal{S}(\hat{\mathrm{G}})\) : IX.8.2. \(\chi_u\) : IX.8.3. \(\mathrm{Z}\mathcal{C}^{\infty}(\mathrm{G}),\mathrm{Z}\mathcal{C}(\mathrm{G})\) : IX.8.4. \(\operatorname {E}_{(t)}\) : IX.9.4.\\
Spin(n,R):IX.3,习题 5.

\specialchapter{术语索引}

这些参考编号分别表示章、节和编号。

绝对单（李代数）：VIII.3.2 伴随（SL(2, k) 的表示）：VIII.1.4 容许（格）：VIII.12.8 单房：IX.5.2 反埃尔米特（自同态）：IX.1.1 反不变：IX.7.4 双序（李代数的）：VIII.12.1 Borel（子代数）：VIII.3.3, VIII.3.5 典则（嘉当子代数、根系、基）：VIII.5.3 典则（sl(2, k) 的对合）：VIII.1.1 嘉当（子代数）：VII.2.1 Casimir 元：IX.7.6 中心（g-模的特征标）：VIII.6.1 中心（函数）：IX.8.3 房：IX.5.2 特征标（g-模的）：VIII.7.7 Chevalley（序）：VIII.12.7 Chevalley（系）：VIII.2.4 经典（可分裂李代数）：VIII.3.2 净（子群）：IX.4, 习题 15 Clebsch-Gordan（公式）：VIII.9.4 紧（李代数）：IX.1.3 紧（复李代数的实形式）：IX.3.1 相容（表示）：VII.3.1, VIII.1.4 条件 (AC)：VII.1.1 共轭：IX.3.1 反变图式：IX.4.9 共变图式：IX.4.9 Coxeter 元：IX.4, 习题 14 支配：IX.7.1 特征值：VII.1.1 特征向量：VII.1.1

初等（自同构）：VII.3.1 嵌入（流形到子集邻域中的）：IX.9.1 例外型（可分裂单李代数）：VIII.3.2 指数：VIII.8.1 面（与抛物子集相联系的）：VIII.3.4 Fitting（模的分解）：VII.1.1 旗：VIII.13.1 Fourier 余变换：IX.8.1 Fourier 变换：IX.8.1 标架：IX.4.10 标架（分裂半单李代数的）：VIII.4.1 带标架的（半单李代数）：VIII.4.1 基本（支配权）：IX.7.1 基本（表示、单 g-模）：VIII.7.2 生成族（由标架定义的）：VIII.4.1 Harish-Chandra（同态）：VIII.6.4 Hermann Weyl（-- 的公式）：VIII.9.1 齐性对称空间：IX.1, 习题 8 不变（多项式函数）：VIII.8.3 对合（李代数）：IX.1, 习题 7 不可约（对合李代数）：IX.1, 习题 7 不可约（齐性对称空间）：IX.1, 习题 8 迷向（旗）：VIII.13.2, VIII.13.3 同构型（最高权 λ 的分支）：VIII.7.2 Jacobson-Morozov（定理）：VIII.11.2 格（Q-向量空间中的）：VIII.12.1 线性管：IX.9.3 极大环面：IX.2.2 极小（权）：VIII.7.3 缓增：IX.8.2 幂零（元素的分量）：VII.1.3 幂零空间：VII.1.1 结点（环面的群）：IX.4.2 结点向量：IX.4.5 序（Q-代数的）：VIII.12.1 正交（李代数）：VIII.13.2 正交（不可约表示）：VIII.7.5 抛物（子代数）：VIII.3.4, VIII.3.5 （正根的）分拆：VIII.9.1 许可（格）：VIII.12.6 正根：IX.7.1 准素（子空间）：VII.1.1 本原（模的元）：VIII.1.2, VIII,6.1 主（幂零元、单元素、sl₂ 三元组）：VIII.11.4 主（轨道）：IX.9.4 主（子群）：IX.4, 习题 18 真（子空间）：VII.1.1 拟极大（迷向旗）：VIII.13.4 R-极值（元）：VIII.7.2 R-饱和（子集）：VIII.7.2 根（子群）：IX.4.7 秩（李代数的）：VII.2.2 秩（李群的）：IX.2.4 速降：IX.8.2 实形式（复李代数的）：IX.3.1 既约（根图式）：IX.4.8 正则（线性表示的元）：VII.4.1 正则（李代数的元）：VII.2.2 正则（李群的元）：VII.4.2 表示环（李代数的）：VIII.7.6 根：IX.4.4 根（(g,h) 的）：VIII.2.2 根图式：IX.4.8 半单（元素的分量）：VII.1.3 半旋量（表示）：VIII.13.4 单（元素）：VIII.11.3 单（根）：IX.7.1 奇异超平面：IX.5.2 sl₂ 三元组：VIII.11.1 旋量（表示）：VIII.13.2, VIII.13.4 分裂（约化李代数）：VIII.2.1 可分裂（李子代数）：VII.5.1, VIII.10 可分裂（嘉当子代数、约化李代数）：VIII.2.1 可分裂包络（gl(V) 中子代数的）：VII.5.2 子群（嘉当子群）：IX.2.2 子群（极大秩子群）：IX.2.4 子群（主子群）：IX.4, 习题 18 辛（李代数）：VIII.13.3 辛（不可约表示）：VIII.7.5 切（多项式映射的线性映射）：VII.App.I.2 紧 Cʳ 收敛拓扑：IX.6.4 挠素数：IX.5, 习题 9 至 11 环面：IX.1.2 横截集：IX.9.3 及 IX.9，习题 16 型（An,Bn,\ldots）：IX.3.3 型（表示的）：IX.App.II.1, IX.App.II.2 向量（群）：IX.1.2 强正则（元）：IX.5.1

权（g-模的）：VII.1.1, VIII.1.2, VIII.6.1

权（表示的）：IX.4.3

权（分裂李代数的权群）：VIII.2.2

Weyl（群）：IX.2.5

Weyl（分裂代数的群）：VIII.2.2

Witt（基）：VIII.13.2, VIII.13.3, VIII.13.4

Zariski（拓扑）：VII.App.I.1